\section{Measure and integration from coherent finite partitions}
\label{sec:measure}

The existing expectation packet already constructs expectations from effective
\(L^1\) simple-law names and supplies compact quadrature and improper-tail
arguments \cite{repoExpectedValue}.  The following construction extends that
interface in three specific directions: several functions live over one
coherent algebra of events; convergence in measure can be converted to
\(L^1\) convergence with an explicit budget; and integration in one variable
extends as a contraction, giving a represented Fubini theorem.  It does not
assert the unrestricted classical dominated convergence theorem in iPRA.
The need for this distinction is substantive: precise second-order versions
of dominated convergence have nontrivial reverse-mathematical strength
\cite{measureADR2012}.  Predicative completion-based integration is an
established approach \cite{measurePZ2024}; the additional restriction here is
to fixed primitive-recursive presentations and their arithmetic certificates.

\subsection{One algebra of events, rather than unrelated couplings}

\begin{definition}[Rational measured Boolean presentation]
Fix a numerical coding of a Boolean algebra \(\mathcal B\), with
primitive-recursive Boolean operations, decidable coded equality, and a
primitive-recursive rational function \(\mu\).  Its displayed laws are
\[
 \mu(0)=0,\qquad \mu(b)\geq0,\qquad \mu(1)=M\in\mathbb Q_{\geq0},
 \qquad b\wedge c=0\Longrightarrow
 \mu(b\vee c)=\mu(b)+\mu(c).
\]
The Boolean identities and these laws must have proofs for the chosen
presentation.  A partition is a finite list \((b_i)_{i<r}\) with pairwise
disjoint entries and join \(1\).  A rational simple function is such a
partition together with rational values \((a_i)_{i<r}\).  Zero cells and
zero-measure cells are allowed.
\end{definition}

All these objects are finite numerical codes.  A concrete continuous example
is the algebra of finite unions of dyadic half-open intervals in \([0,1)\),
with rational length.  Represent an event by a level and its bit vector,
raise two levels to their maximum, and repeat the appropriate bits before
applying Boolean operations.  Length is the number of selected cells divided
by the number of cells.  This gives primitive-recursive operations and finite
proofs of the measure laws; it does not introduce an arbitrary measurable-set
parameter into first-order arithmetic.

Given partitions \((b_i)\) and \((c_j)\), use the partition
\((b_i\wedge c_j)_{i,j}\).  It is a common refinement, and a third partition
refines it by another meet.  Associativity of meet supplies coherent joint
refinements for any finite list of simple functions.  This coherence is an
extra condition if one starts instead from pairwise couplings of probability
laws: unrelated pairwise couplings alone do not specify a common joint
probability space for three or more variables.

For dyadic presentations, implement a batch of comparisons on the finest
requested grid once.  The common grid has \(2^{\max m_i}\) cells, whereas
repeated formal products of partitions can introduce many redundant empty
cells.  Retain the common-grid code and its mass array for every operation
in the batch.  This is a concrete efficiency improvement with exactly the
same refinement proof; no polynomial bound is asserted for arbitrary Boolean
presentations.

Define
\[
 I(s)=\sum_{i<r}a_i\mu(b_i),\qquad
 \|s\|_1=\sum_{i<r}|a_i|\mu(b_i),\qquad d_1(s,t)=\|s-t\|_1,
\]
where arithmetic on two simple functions uses their common refinement.

\begin{lemma}[Finite integration calculus]\label{lem:measure-finite}
Refinement preserves \(I\) and \(\|\cdot\|_1\).  On simple functions,
\[
 |I(s)-I(t)|\leq\|s-t\|_1,\qquad
 \bigl|\|s\|_1-\|t\|_1\bigr|\leq\|s-t\|_1,
 \qquad \||s|-|t|\|_1\leq\|s-t\|_1.
\]
Integration is linear and positive; \(d_1\) is a pseudometric.  The maps
\(s\mapsto s^+=\max(s,0)\) and
\(s\mapsto T_Ks=\max(-K,\min(K,s))\), for rational \(K\geq0\),
are nonexpansive.  Multiplication by a simple function bounded by a supplied
rational \(C\) is \(C\)-Lipschitz in \(d_1\).
\end{lemma}

\begin{proof}
If \(b_i\) is subdivided into disjoint cells \(d_{ij}\), finite additivity
gives \(\mu(b_i)=\sum_j\mu(d_{ij})\).  Replacing
\(a_i\mu(b_i)\) by \(\sum_j a_i\mu(d_{ij})\) preserves the integral;
the same identity with \(|a_i|\) preserves the norm.  Thus all comparisons
may be made on one partition.  On its \(i\)-th cell, use
\(|a_i+b_i|\leq|a_i|+|b_i|\),
\(\bigl||a_i|-|b_i|\bigr|\leq|a_i-b_i|\), and the elementary rational
inequalities for maximum, minimum, and multiplication by a bounded number.
Multiply by the nonnegative cell mass and sum.  The sum of nonnegative
terms is nonnegative, proving positivity as well.  Nothing requires deciding
whether a completed real number is zero or positive.

For an object-language proof, clear the positive rational denominators.
Induct on the length of a partial sum using the corresponding inequality
between its two accumulators.  The partition laws and all cell comparisons
are primitive-recursive Boolean checks, and each accumulator inequality is
quantifier-free.  Universal hypotheses about a presentation are instantiated
at the cells actually used; they are not induction predicates.
\end{proof}

\subsection{Completion, order, and the exact adequacy claim}

An \(L^1\) name is a fixed primitive-recursive family \(s_k\) of simple
functions with proofs of
\begin{equation}
 j\geq k\quad\Longrightarrow\quad
 \|s_j-s_k\|_1\leq2^{-k}.
 \label{eq:measure-regular}
\end{equation}
This is a convenient one-sided convention.  From a name with the symmetric
regularity bound \(\|s_j-s_k\|_1\leq2^{-j}+2^{-k}\), as used by the
common completion interface, set \(t_k=s_{k+1}\).  For \(j\geq k\),
\[
 \|t_j-t_k\|_1\leq2^{-(j+1)}+2^{-(k+1)}\leq2^{-k}.
\]
The shift denotes the same represented object.  Conversely, a one-sided
name satisfies symmetric regularity by ordering its two numerical indices.
Its limit notation \(f\) means this represented completion element.
Equality, inequalities, and norms are interpreted by the completion's
finite-precision relations.  In particular, \(\|f-s_k\|_1\leq2^{-k}\)
abbreviates the corresponding family of rational observations.  It does
not supply a pointwise evaluator for an arbitrary integrable function.
Here ``completion'' means the collection of these certified names, embedded
in the external full metric completion.  Closure is asserted for supplied
jointly primitive-recursive Cauchy families with explicit moduli; the
primitive-recursive names are not asserted to form a classically complete
Banach space under arbitrary Cauchy sequences.

\begin{theorem}[Represented integration and lattice operations]
\label{thm:measure-completion}
The sequences \(I(s_k)\) and \(\|s_k\|_1\) are regular rational names.
They define a well-defined positive linear integral and norm on the
represented completion.  Absolute value, positive part, and clipping extend
nonexpansively.  Addition uses \(s_{k+1}+t_{k+1}\); multiplication by a fixed
rational \(c\) uses \(cs_{k+b}\), where \(|c|\leq2^b\).  These choices
give explicit regular output names.  More generally let \(b_k\) be simple
functions with supplied bounds
\(\|b_k\|_\infty\leq C\) and
\(\|b_j-b_k\|_\infty\leq2^{-k}\) for \(j\geq k\), where the finite
supremum is the maximum of the absolute cell values.  Then multiplication
by the represented bounded multiplier is defined by the products
\(s_{k+h}b_{k+h}\), where
\(C+U\leq2^h\) and \(U=\|s_0\|_1+1\).
\end{theorem}

\begin{proof}
Apply Lemma~\ref{lem:measure-finite} at two indices in
\eqref{eq:measure-regular}.  This proves regularity of the scalar output
names and of every nonexpansive operation.  If two input names have zero
represented distance, the same inequalities force zero distance of their
outputs.  Addition loses the sum of two errors, which the shift by one
compensates; scalar multiplication loses \(|c|\), which the shift by \(b\)
compensates.  For the last assertion, \(\|s_k\|_1\leq U\) and
\[
 \|s_jb_j-s_kb_k\|_1
 \leq C\|s_j-s_k\|_1+U\|b_j-b_k\|_\infty
 \leq(C+U)2^{-k};
\]
the displayed shift makes the output regular.  Its extensionality follows
from the same estimate.  Linearity follows from the exact identity at any common finite
level, paying these already specified approximation errors.  Positivity is
defined by equality with the positive part, or equivalently zero norm of the
negative part, and the finite positive-integral law passes through the same
observations.  Every requested conclusion uses finitely many indices chosen
by these explicit schedules.
\end{proof}

For clarity, the completed order is
\(f\leq g\) when \(\|(f-g)^+\|_1=0\).  This is order on classes modulo
zero \(L^1\) distance.  It neither asserts nor tests pointwise order at every
point.  The usual lattice identities follow from their exact finite
identities and the nonexpansive estimates.

The events themselves have a useful completion.  For coded events,
\(d_\mu(b,c)=\mu(b\mathbin{\triangle}c)
=\|\mathbf1_b-\mathbf1_c\|_1\).  A represented event is a regular name
of indicators for this pseudometric.  Complement is an isometry, while
\[
 d_\mu(b\wedge c,b'\wedge c'),\quad
 d_\mu(b\vee c,b'\vee c')
 \ \leq d_\mu(b,b')+d_\mu(c,c').
\]
Indeed, changing a meet or union can only change a cell where at least one
of its two input indicators changes; this is a finite Boolean identity and
weighted inequality.  Thus complement extends directly and meet and union
extend with the input shift \(k\mapsto k+1\).  Their Boolean laws pass
through finite comparisons, and the identity
\(\mu(b\vee c)+\mu(b\wedge c)=\mu(b)+\mu(c)\) survives completion.
This constructs a represented measure algebra, with exact equality understood
modulo zero measure.  It uses bounded indicator products, and does not assert
unrestricted multiplication of \(L^1\) classes.

There is a precise external interpretation.  Suppose the coded algebra is
realized as measurable sets of an ordinary finite measure space, its rational
measure agrees with the ambient measure, and an ordinary integrable function
\(F\) comes with a primitive-recursive family \(s_k\) and proved estimates
\(\int|F-s_k|\leq2^{-k}\).  These estimates give symmetric regularity,
so use \(t_k=s_{k+1}\) to obtain \eqref{eq:measure-regular}.
The represented integral agrees with \(\int F\), since
\[
 \left|I(t_k)-\int F\right|\leq\int|s_{k+1}-F|\leq2^{-(k+1)}.
\]
If \(u_k\) is another such approximation family for \(F\), the triangle
inequality gives \(\|s_{k+1}-u_{k+1}\|_1\leq2^{-k}\), so both shifted
names denote the same completion
element.  Conversely, in ordinary measure theory a name satisfying
\eqref{eq:measure-regular} has an \(L^1\) limit by completeness.  This last
sentence explains the standard semantics; it is not used as an additional
iPRA axiom or an algorithm producing a pointwise representative.  We claim
adequacy for functions with the supplied approximation data, not that every
classical Lebesgue-integrable function possesses such data.

\begin{proposition}[Countable unions with an explicit tail]
\label{prop:measure-union}
Let \(b_j\) be a fixed primitive-recursive event sequence.  Suppose a
primitive-recursive \(C\) satisfies
\[
 L\geq N\geq C(k)\quad\Longrightarrow\quad
 \sum_{N\leq j<L}\mu(b_j)\leq2^{-k}.
\]
The indicators of the finite unions \(\bigvee_{j<N}b_j\), reindexed by a
monotone version of \(C\), define a represented event.  If the events are
pairwise disjoint, its measure is the represented sum \(\sum_j\mu(b_j)\).
\end{proposition}

\begin{proof}
The difference of two union indicators is supported on the union of the
intervening events and is bounded by the sum of their indicators.  Therefore
its \(L^1\) norm is at most the displayed tail sum.  For disjoint events,
finite additivity identifies the measure of every finite union with the
finite sum, and the scalar names agree.  The tail is part of the input; no
unrestricted countable additivity or automatic modulus extraction is inferred.
\end{proof}

\subsection{A quantitative Vitali theorem and a dominated special case}

For \(K\geq0\), write
\(R_K(t)=(|t|-K)^+\).  The rational functions \(T_K\) and \(R_K\) are
1-Lipschitz, and \(|t-T_K(t)|=R_K(t)\).

\begin{lemma}[Finite small-exception estimate]\label{lem:measure-vitali-finite}
Let \(u,v\) be rational simple functions on a common partition, let
\(\delta>0\), and let \(B\) be the union of cells where
\(|u-v|>\delta\).  If \(\mu(B)\leq\eta\), then
\begin{equation}
 \|u-v\|_1\leq I(R_Ku)+I(R_Kv)+\delta M+2K\eta.
 \label{eq:measure-vitali-finite}
\end{equation}
\end{lemma}

\begin{proof}
The triangle inequality splits \(u-v\) into its two clipping remainders and
\(T_Ku-T_Kv\).  Off \(B\), the last difference is at most \(\delta\)
because clipping is nonexpansive.  On \(B\), it is at most \(2K\).
Integrating the finite cellwise bound
\[
 |u-v|\leq R_Ku+R_Kv+\delta+2K\mathbf1_B
\]
proves the result.  In arithmetic this is a Boolean branch at a comparison of
rational numbers, followed by the weighted partial-sum invariant of
Lemma~\ref{lem:measure-finite}.  There is no comparison of completed real
values at a point.
\end{proof}

\begin{lemma}[An \(L^1\) majorant supplies its own tail modulus]
\label{lem:measure-majorant-tail}
Let \(g\geq0\) have an \(L^1\) name \(g_r\).  If a rational \(K\)
bounds every value of \(g_r\), in absolute value, then
\[
 I((g-K)^+)\leq\|g-g_r\|_1\leq2^{-r}.
\]
Consequently, if \(|f_n|\leq g\) and \(|f|\leq g\) in the completed
lattice, then both \(I(R_Kf_n)\) and \(I(R_Kf)\) are at most \(2^{-r}\),
uniformly in \(n\).
\end{lemma}

\begin{proof}
On a common finite partition, the scalar inequality
\((a-K)^+\leq|a-b|\) holds whenever \(|b|\leq K\).
It passes to represented names by nonexpansiveness and positivity.
The map \(a\mapsto(a-K)^+\) preserves order, so domination by \(g\)
gives the two further bounds.  At required tail error \(2^{-r}\), one
therefore evaluates the single simple approximant \(g_r\) and takes the
maximum of its finitely many absolute rational values.  This maximum is a
primitive-recursive operation; a separate amplitude-tail oracle is unnecessary.
\end{proof}

\begin{theorem}[Vitali transfer with an explicit precision budget]
\label{thm:measure-vitali}
Let \(f_n\) and \(f\) be uniformly specified \(L^1\) names.  Fix
\(\varepsilon=2^{-k}\).  Suppose primitive-recursive constructions provide
rationals \(K\geq1,\delta>0,\eta\geq0\), a precision \(r\), and a
cutoff \(N(k)\), such that for every \(n\geq N(k)\):
\begin{enumerate}
 \item the chosen simple approximants \(u\) to \(f_n\) and \(v\) to
 \(f\) have \(L^1\) errors at most \(\sigma\);
 \item \(I(R_Kf_n),I(R_Kf)\leq\tau\);
 \item the finite event \(B=\{|u-v|>\delta\}\) has measure at most
 \(\eta\).
\end{enumerate}
Then
\begin{equation}
 \|f_n-f\|_1\leq4\sigma+2\tau+\delta M+2K\eta.
 \label{eq:measure-vitali-budget}
\end{equation}
In particular, the choices
\begin{equation}
 \sigma\leq\frac{\varepsilon}{16},\quad
 \tau\leq\frac{\varepsilon}{8},\quad
 \delta=\frac{\varepsilon}{4(M+1)},\quad
 \eta=\frac{\varepsilon}{8(K+1)}
 \label{eq:measure-vitali-schedule}
\end{equation}
give \(\|f_n-f\|_1\leq\varepsilon\), and hence
\(|I(f_n)-I(f)|\leq\varepsilon\).
\end{theorem}

\begin{proof}
The outer approximation loses \(2\sigma\):
\(\|f_n-f\|_1\leq2\sigma+\|u-v\|_1\).
Since \(R_K\) is nonexpansive,
\(I(R_Ku)\leq I(R_Kf_n)+\sigma\leq\tau+\sigma\), and likewise
\(I(R_Kv)\leq\tau+\sigma\).  Substitute both inequalities into
\eqref{eq:measure-vitali-finite}.  This proves
\eqref{eq:measure-vitali-budget}.  In
\eqref{eq:measure-vitali-schedule}, each of its four summands is at most
\(\varepsilon/4\).  The scalar integral estimate follows from
Lemma~\ref{lem:measure-finite} and completion.
\end{proof}

For the dominated special case, take
\(K=1+\max|g_{k+3}|\) in Lemma~\ref{lem:measure-majorant-tail}.
This gives \(\tau\leq2^{-(k+3)}\).  Choose \(r\geq k+4\) for the
simple approximants, and supply the displayed finite small-bad-set
certificate with its cutoff \(N(k)\).  The resulting theorem is a dominated
convergence theorem with a quantitative convergence-in-measure witness.
It does not obtain that witness from bare pointwise convergence.

This interface has the expected standard interpretation.  Take \(\eta>0\),
as in \eqref{eq:measure-vitali-schedule}.  In an ambient measure space,
if a modulus already guarantees
\(\mu(|f_n-f|>\delta/2)\leq\eta/2\), and simple approximants have
errors at most \(\sigma\leq\delta\eta/8\), then Markov's finite-sum
inequality, extended by the integral estimate, gives
\[
 \mu(|u-v|>\delta)
 \leq\mu(|f_n-f|>\delta/2)
    +\frac{2}{\delta}\bigl(\|u-f_n\|_1+\|v-f\|_1\bigr)
 \leq\eta.
\]
Thus the required approximant precision may be chosen to satisfy both
\(2^{-r}\leq\varepsilon/16\) and
\(2^{-r}\leq\delta\eta/8\).  Finding such an \(r\) is bounded search
on positive rational numerators and denominators.  This explanation uses
ordinary level-set semantics only to connect a conventional input modulus
to the finite certificate.  The arithmetic transfer theorem itself uses
only the explicitly computed finite event \(B\).

\paragraph{An unbounded moving-spike example.}
In \([0,1)\), let
\[
 A_j=[4^{-(j+1)},4^{-j}),\qquad
 g=\sum_{j\geq0}2^j\mathbf1_{A_j},\qquad
 f_n=2^n\mathbf1_{A_n}.
\]
The endpoint convention affects only finitely or countably many null points
in the ordinary interpretation.  Exactly,
\[
 \mu(A_j)=\frac{3}{4^{j+1}},\qquad
 I(g)=\frac32,\qquad
 \left\|g-\sum_{j<J}2^j\mathbf1_{A_j}\right\|_1
       =\frac32\,2^{-J},\qquad
 \|f_n\|_1=\frac34\,2^{-n}.
\]
These are finite geometric-sum identities with their explicit tails.
The functions \(f_n\) converge pointwise to zero in the usual realization,
are dominated by the single represented integrable \(g\), and have growing
supremum \(2^n\); uniform convergence and a common pointwise constant bound
are unavailable.  Their finite bad sets are \(A_n\), so the hypotheses of
Theorem~\ref{thm:measure-vitali} have direct primitive-recursive witnesses.
There is also a useful boundary check: \(g\in L^1\), but the integral of
the square of its first \(J\) pieces is \(3J/4\).  Therefore the framework
must not declare arbitrary multiplication \(L^1\times L^1\to L^1\).
The bounded-multiplier hypothesis above has mathematical content.

\paragraph{Explicit exhaustions of infinite-measure spaces.}
The same norm construction can start from a ring of finite-measure events
and simple functions of finite support; only the small-exception estimate
needs a total-mass bound.  For that estimate, restrict to a supplied finite
event \(E_k\) of mass \(M_k\).  If the input additionally proves
\[
 \|f_n\mathbf1_{E_k^c}\|_1+
 \|f\mathbf1_{E_k^c}\|_1\leq\rho_k
 \qquad(n\geq N(k)),
\]
then apply the finite theorem on \(E_k\), with \(M_k\) in place of \(M\),
and add \(\rho_k\) to the output error.  Here \(\mathbf1_{E_k^c}f\)
means \(f-\mathbf1_{E_k}f\), so no finite mass is assigned to the complement.
An explicit exhaustion and these uniform spatial-tail bounds therefore
cover improper integrals on \(\mathbb R^d\).  The proof is the triangle
inequality after this split; neither abstract sigma-finiteness nor a bare
assertion of integrability provides the needed schedules.

\subsection{Fubini as extension of a finite contraction}

Take two rational measured Boolean presentations \((\mathcal B_X,\mu)\)
and \((\mathcal B_Y,\nu)\).  A rectangular simple function is a rational
matrix \((a_{ij})\) on finite partitions \((b_i)\) and \((c_j)\), with
product-cell mass \(\mu(b_i)\nu(c_j)\).  Common row and column refinements
make the following operations independent of presentation:
\[
 I_{XY}(a)=\sum_{i,j}a_{ij}\mu(b_i)\nu(c_j),\qquad
 (J_Ya)_i=\sum_j a_{ij}\nu(c_j),\qquad
 (J_Xa)_j=\sum_i a_{ij}\mu(b_i).
\]

\begin{theorem}[Represented Fubini]\label{thm:measure-fubini}
The operators \(J_Y\) and \(J_X\) extend to the \(L^1\) completion of
rectangular simple functions.  They are linear contractions, and
\[
 I_X(J_Yf)=I_{XY}(f)=I_Y(J_Xf).
\]
If \(a_k\) is a regular rectangular name, the regular output names are
\(J_Ya_k\) and \(J_Xa_k\): no additional precision shift is necessary.
\end{theorem}

\begin{proof}
At the finite level,
\[
 \|J_Ya\|_{L^1(X)}
 =\sum_i\mu(b_i)\left|\sum_j a_{ij}\nu(c_j)\right|
 \leq\sum_{i,j}|a_{ij}|\mu(b_i)\nu(c_j)
 =\|a\|_{L^1(XY)}.
\]
Apply this to a difference of two matrices on a common refinement.  It proves
both regularity of the output and independence of equivalent names.
Finite linearity proves linearity after completion.  Interchanging the two
finite sums is an exact rational identity, so the three scalar output names
agree at each finite level.  Thus their represented limits agree as well.
For iPRA, the contraction is an iterated weighted-sum inequality and the
interchange identity has a quantifier-free accumulator proof.  The existing
object-level finite-grid module already illustrates the required separation
between cell hypotheses and the two accumulated bounds \cite{repoGridIntegral}.
\end{proof}

This conclusion supplies the iterated integral as an \(L^1\) class on the
remaining variable.  It does not assert that every point has an integrable
section, nor select representatives of all classes.  In standard measure
semantics, the represented operator agrees with ordinary partial integration
on its domain of functions with the supplied rectangular \(L^1\) names.
The finite theorem and its completion proof remain valid without building
an entire product sigma-algebra inside iPRA.  Complex-valued integration is
obtained by two real copies, using the rational norm
\(|\operatorname{Re}z|+|\operatorname{Im}z|\); this gives the same
integrability notion as the ordinary complex norm, with explicit equivalent
norm bounds.

\subsection{Lean and object-language proof obligations}

The following are the additional definitions needed beyond the common
completion interface.  A \texttt{MeasuredBooleanCode} stores numerical
event operations and rational mass.  A \texttt{PartitionCode} stores a
finite event array and the Boolean validity checker; a
\texttt{SimpleCode} adds a rational value array.  Define
\texttt{refine} by pairwise meet, \texttt{integral} and \texttt{l1} by
primitive-recursive folds, \texttt{clip} and \texttt{tail} by rational
minimum/maximum, and \texttt{badSet} by a finite Boolean union.
\texttt{RectSimpleCode} adds a second partition and a rectangular array.
These finite definitions have no quotient, choice operation, or quantified
set of measurable functions.

The proof order is: refinement invariance; scalar clipping identities;
weighted integral and lattice inequalities; the finite bad-set estimate;
the four-term budget; and row/column contraction and sum interchange.
Only then instantiate the completion extension theorem.  At the formal
arithmetic level, export the denominator-cleared finite statements as
\texttt{Prf PRAAx} derivations, with partition validity and finite comparison
checks as hypotheses.  Supply fixed primitive-recursive terms for the input
names, majorant maximum, bad-set generator, and cutoff; prove their contracts
separately.  The concluding named-integral equalities are families of
finite-precision consequences of those derivations.  A Lean structure
quantifying over arbitrary functions is acceptable as metatheoretic notation,
but cannot itself replace these syntactic proof obligations.

The accompanying exact-rational program checks the finite estimate on
enumerated small weighted examples, verifies the rectangular contraction and
sum identity, and reproduces the moving-spike formulas.  These calculations
test the implementation of the displayed finite formulas.  They are not
evidence that arbitrary dominated convergence, Caratheodory extension,
Radon--Nikodym differentiation, or a completed iPRA formalization has been
obtained.  General metastable convergence theorems also require separate
complexity auditing: the recursion used in a general quantitative theorem
need not be primitive recursive \cite{measureADRMeta2012}.
